\chapter{Chapter09 — Friedgut}
%%% generated by the blueprint pipeline from FABL.Chapter09.Friedgut
\section{Overview}

Declarations of the Lean module \texttt{FABL.Chapter09.Friedgut}.

\section{Declarations}

\begin{definition}[Influential coordinates]
\lean{FABL.influentialCoordinates}
\statementsource{Analysis-of-Boolean-Functions-by-Ryan-ODonnell}{
  pdf-ffbe9c7b3306-p046-b003,
  pdf-ffbe9c7b3306-p047-b005,
  pdf-ffbe9c7b3306-p215-b006,
  pdf-ffbe9c7b3306-p215-b007,
  pdf-ffbe9c7b3306-p268-b003
}
\label{FABL.influentialCoordinates}
\leanok
\uses{FABL.BooleanFunction, FABL.booleanInfluence}
For a Boolean function $f$ on the $n$-dimensional Boolean cube and a real threshold
$\tau$, the influential coordinates are the indices $i\in\{1,\ldots,n\}$ whose Boolean
influence is at least $\tau$.
\end{definition}

\begin{theorem}[Membership in the influential coordinates]
\lean{FABL.mem_influentialCoordinates}
\label{FABL.mem_influentialCoordinates}
\leanok
\uses{FABL.BooleanFunction, FABL.booleanInfluence, FABL.influentialCoordinates}
\difficulty{1}
Let $n$ be a natural number, let $f$ be a Boolean function on the $n$-dimensional
Boolean cube, let $\tau\in\mathbb{R}$, and let $i$ be a coordinate. Then $i$ is among
the coordinates whose influence under $f$ is at least $\tau$ if and only if $\tau$ is
at most the influence of $i$ under $f$.
\end{theorem}

\begin{theorem}[Not mem influential coordinates]
\lean{FABL.not_mem_influentialCoordinates}
\statementsource{Analysis-of-Boolean-Functions-by-Ryan-ODonnell}{pdf-ffbe9c7b3306-p268-b003}
\label{FABL.not_mem_influentialCoordinates}
\leanok
\uses{FABL.BooleanFunction, FABL.booleanInfluence, FABL.influentialCoordinates}
\difficulty{1}
% dataset-pipeline: stub -- prose pending build_informal_statements.py
Statement pending regeneration.
\end{theorem}

\begin{theorem}[Influence bound for influential coordinates]
\lean{FABL.card_influentialCoordinates_mul_le_totalInfluence}
\label{FABL.card_influentialCoordinates_mul_le_totalInfluence}
\leanok
\uses{FABL.BooleanFunction, FABL.BooleanFunction.toReal, FABL.booleanInfluence_eq_influence_toReal, FABL.influence, FABL.influence_nonneg, FABL.influentialCoordinates, FABL.mem_influentialCoordinates, FABL.totalInfluence}
\difficulty{4}
Let $n$ be a natural number, let $f:\{-1,1\}^n\to\{-1,1\}$ be a Boolean function,
and let $\tau$ be any real number. If $S$ is the set of coordinates whose influence is
at least $\tau$, then $|S|\tau\leq I$, where $I$ is the total influence of the
associated real-valued function.
\end{theorem}

\begin{theorem}[Influential coordinate count bound]
\lean{FABL.card_influentialCoordinates_le_div}
\proofsource{Analysis-of-Boolean-Functions-by-Ryan-ODonnell}{
  pdf-ffbe9c7b3306-p215-b007,
  pdf-ffbe9c7b3306-p268-b003,
  pdf-ffbe9c7b3306-p268-b004
}
\label{FABL.card_influentialCoordinates_le_div}
\leanok
\uses{FABL.BooleanFunction, FABL.BooleanFunction.toReal, FABL.card_influentialCoordinates_mul_le_totalInfluence, FABL.influentialCoordinates, FABL.totalInfluence}
\difficulty{1}
Let $n$ be a natural number and let $f:\{-1,1\}^n\to\{-1,1\}$ be a Boolean
function. For a real threshold $\tau>0$, let $S$ be the coordinates whose influence is
at least $\tau$. Then $|S|\leq I/\tau$, where $I$ is the total influence of the
associated real-valued function.
\end{theorem}

\begin{definition}[Low degree coordinate mass]
\lean{FABL.lowDegreeCoordinateMass}
\statementsource{Analysis-of-Boolean-Functions-by-Ryan-ODonnell}{
  pdf-ffbe9c7b3306-p028-b002,
  pdf-ffbe9c7b3306-p269-b001
}
\label{FABL.lowDegreeCoordinateMass}
\leanok
\uses{FABL.BooleanFunction, FABL.BooleanFunction.toReal, FABL.fourierCoeff}
For a Boolean function $f$ on the $n$-dimensional Boolean cube, a coordinate
$i\in\{1,\ldots,n\}$, and a nonnegative integer $k$, the low degree mass associated with
$i$ is the sum of $\widehat f(S)^2$ over all subsets $S\subseteq\{1,\ldots,n\}$ that
contain $i$ and have cardinality at most $k$.
\end{definition}

\begin{theorem}[Low-degree mass and stable influence at one third]
\lean{FABL.lowDegreeCoordinateMass_le_three_pow_mul_stableInfluence}
\proofsource{Analysis-of-Boolean-Functions-by-Ryan-ODonnell}{
  pdf-ffbe9c7b3306-p057-b001,
  pdf-ffbe9c7b3306-p269-b001
}
\label{FABL.lowDegreeCoordinateMass_le_three_pow_mul_stableInfluence}
\leanok
\uses{FABL.BooleanFunction, FABL.BooleanFunction.toReal, FABL.fourierCoeff, FABL.lowDegreeCoordinateMass, FABL.stableInfluence}
\difficulty{5}
Let $f$ be a sign-valued Boolean function on the $n$-dimensional Boolean cube, let $i$
be a coordinate, and let $k$ be a natural number. The low-degree coordinate mass of $f$
at $i$ through degree $k$ is at most $3^k$ times the stable influence of the real-valued
version of $f$ at $i$ with parameter $1/3$.
\end{theorem}

\begin{definition}[Low degree mass outside influential coordinates]
\lean{FABL.lowDegreeOutsideInfluentialMass}
\statementsource{Analysis-of-Boolean-Functions-by-Ryan-ODonnell}{
  pdf-ffbe9c7b3306-p069-b003,
  pdf-ffbe9c7b3306-p269-b001
}
\label{FABL.lowDegreeOutsideInfluentialMass}
\leanok
\uses{FABL.BooleanFunction, FABL.BooleanFunction.toReal, FABL.fourierCoeff, FABL.influentialCoordinates}
For a Boolean function $f$ on the $n$-dimensional Boolean cube, a real threshold $\tau$,
and a nonnegative integer $k$, this mass is the sum of $\widehat f(S)^2$ over all
subsets $S\subseteq\{1,\ldots,n\}$ of size at most $k$ that are not contained in the set
of coordinates whose Boolean influence is at least $\tau$.
\end{definition}

\begin{theorem}[Coordinate bound for low degree mass]
\lean{FABL.lowDegreeOutsideInfluentialMass_le_sum_coordinateMass}
\label{FABL.lowDegreeOutsideInfluentialMass_le_sum_coordinateMass}
\leanok
\uses{FABL.BooleanFunction, FABL.BooleanFunction.toReal, FABL.booleanInfluence, FABL.fourierCoeff, FABL.influentialCoordinates, FABL.lowDegreeCoordinateMass, FABL.lowDegreeOutsideInfluentialMass, FABL.not_mem_influentialCoordinates}
\difficulty{6}
Let $f$ be a function from the $n$-dimensional Boolean cube to $\{-1,1\}$, let
$\tau\in\mathbb R$, and let $k\in\mathbb N$. The sum of the squared Fourier coefficients
over sets $S$ with $|S|\le k$ that are not contained in the coordinates whose influence
is at least $\tau$ is at most the sum, over all coordinates whose influence is less than
$\tau$, of the squared Fourier coefficients over sets $S$ with $|S|\le k$ that contain
that coordinate.
\end{theorem}

\begin{theorem}[Low degree coordinate mass le three pow mul sqrt threshold mul influence]
\lean{FABL.lowDegreeCoordinateMass_le_three_pow_mul_sqrt_threshold_mul_influence}
\statementsource{Analysis-of-Boolean-Functions-by-Ryan-ODonnell}{
  pdf-ffbe9c7b3306-p057-b001,
  pdf-ffbe9c7b3306-p057-b002,
  pdf-ffbe9c7b3306-p268-b004
}
\label{FABL.lowDegreeCoordinateMass_le_three_pow_mul_sqrt_threshold_mul_influence}
\leanok
\uses{FABL.BooleanFunction, FABL.BooleanFunction.toReal, FABL.booleanInfluence, FABL.booleanInfluence_eq_influence_toReal, FABL.influence_nonneg, FABL.influentialCoordinates, FABL.lowDegreeCoordinateMass, FABL.lowDegreeCoordinateMass_le_three_pow_mul_stableInfluence, FABL.not_mem_influentialCoordinates, FABL.stableInfluence, FABL.stableInfluence_oneThird_le_booleanInfluence_mul_sqrt}
\difficulty{4}
% dataset-pipeline: stub -- prose pending build_informal_statements.py
Statement pending regeneration.
\end{theorem}

\begin{theorem}[Low degree outside influential mass le]
\lean{FABL.lowDegreeOutsideInfluentialMass_le}
\statementsource{Analysis-of-Boolean-Functions-by-Ryan-ODonnell}{
  pdf-ffbe9c7b3306-p268-b004,
  pdf-ffbe9c7b3306-p269-b001,
  pdf-ffbe9c7b3306-p057-b001,
  pdf-ffbe9c7b3306-p057-b002
}
\label{FABL.lowDegreeOutsideInfluentialMass_le}
\leanok
\uses{FABL.BooleanFunction, FABL.BooleanFunction.toReal, FABL.booleanInfluence, FABL.booleanInfluence_eq_influence_toReal, FABL.influence_nonneg, FABL.influentialCoordinates, FABL.lowDegreeCoordinateMass, FABL.lowDegreeCoordinateMass_le_three_pow_mul_sqrt_threshold_mul_influence, FABL.lowDegreeOutsideInfluentialMass, FABL.lowDegreeOutsideInfluentialMass_le_sum_coordinateMass, FABL.totalInfluence}
\difficulty{4}
% dataset-pipeline: stub -- prose pending build_informal_statements.py
Statement pending regeneration.
\end{theorem}

\begin{definition}[Fourier influence threshold]
\lean{FABL.theorem9_28Threshold}
\statementsource{Analysis-of-Boolean-Functions-by-Ryan-ODonnell}{pdf-ffbe9c7b3306-p268-b003}
\label{FABL.theorem9_28Threshold}
\leanok
\uses{FABL.BooleanFunction, FABL.BooleanFunction.toReal, FABL.totalInfluence}
For a Boolean function $f$ on the $n$-dimensional Boolean cube, a real number
$\varepsilon$, and a nonnegative integer $k$, define the threshold
$\varepsilon^2/I(f)^2\cdot 9^{-k}$, where $I(f)$ is the total influence of $f$.
\end{definition}

\begin{definition}[Theorem 9.28 family]
\lean{FABL.theorem9_28Family}
\statementsource{Analysis-of-Boolean-Functions-by-Ryan-ODonnell}{pdf-ffbe9c7b3306-p268-b003}
\label{FABL.theorem9_28Family}
\leanok
\uses{FABL.BooleanFunction, FABL.influentialCoordinates, FABL.theorem9_28Threshold}
For a Boolean function $f$ on the $n$-dimensional Boolean cube, a real number
$\varepsilon$, and a nonnegative integer $k$, define a family of subsets of
$\{1,\ldots,n\}$ consisting of those contained in the coordinates whose Boolean
influence is at least $\varepsilon^2/I(f)^2\cdot 9^{-k}$, together with all subsets of
cardinality greater than $k$, where $I(f)$ is the total influence of $f$.
\end{definition}

\begin{theorem}[Positivity of the influence threshold]
\lean{FABL.theorem9_28Threshold_pos}
\label{FABL.theorem9_28Threshold_pos}
\leanok
\uses{FABL.BooleanFunction, FABL.BooleanFunction.toReal, FABL.theorem9_28Threshold, FABL.totalInfluence}
\difficulty{1}
Let $f:\{-1,1\}^n\to\{-1,1\}$, let $\varepsilon>0$, and let $k$ be a nonnegative
integer. If the total influence of $f$ is positive, then
$\dfrac{\varepsilon^2}{I(f)^2}\left(\dfrac{1}{9}\right)^k>0$, where $I(f)$ denotes the
total influence of $f$.
\end{theorem}

\begin{theorem}[Square root of the influence threshold]
\lean{FABL.sqrt_theorem9_28Threshold}
\label{FABL.sqrt_theorem9_28Threshold}
\leanok
\uses{FABL.BooleanFunction, FABL.BooleanFunction.toReal, FABL.theorem9_28Threshold, FABL.totalInfluence}
\difficulty{4}
Let $f$ be a Boolean function on the $n$-dimensional sign cube, let $\varepsilon>0$, and
let $k$ be a natural number. If the total influence of the real-valued version of $f$ is
positive, then the square root of its threshold $\varepsilon^2/I(f)^2\cdot 9^{-k}$
equals $\varepsilon/I(f)\cdot(1/3)^k$, where $I(f)$ denotes that total influence.
\end{theorem}

\begin{theorem}[Bound on influential coordinates]
\lean{FABL.card_theorem9_28_influentialCoordinates_le}
\proofsource{Analysis-of-Boolean-Functions-by-Ryan-ODonnell}{pdf-ffbe9c7b3306-p268-b003}
\label{FABL.card_theorem9_28_influentialCoordinates_le}
\leanok
\uses{FABL.BooleanFunction, FABL.BooleanFunction.toReal, FABL.card_influentialCoordinates_le_div, FABL.influentialCoordinates, FABL.theorem9_28Threshold, FABL.theorem9_28Threshold_pos, FABL.totalInfluence}
\difficulty{3}
Let $n$ and $k$ be natural numbers, and let $f:\{-1,1\}^n\to\{-1,1\}$ be a Boolean
function. Write $I$ for its total influence, and suppose $I>0$. For $\epsilon>0$, the
number of coordinates whose probability of being pivotal for $f$ is at least
$\varepsilon^2 I^{-2}9^{-k}$ is at most $I^3\varepsilon^{-2}9^k$.
\end{theorem}

\begin{theorem}[Exterior weight and the threshold family]
\lean{FABL.fourierWeightOutside_theorem9_28Family}
\label{FABL.fourierWeightOutside_theorem9_28Family}
\leanok
\uses{FABL.BooleanFunction, FABL.BooleanFunction.toReal, FABL.fourierWeight, FABL.fourierWeightOutside, FABL.lowDegreeOutsideInfluentialMass, FABL.theorem9_28Family, FABL.theorem9_28Threshold}
\difficulty{3}
Let $n$ and $k$ be natural numbers, let $f$ be a sign-valued function on the
$n$-dimensional sign cube, and let $\epsilon\in\mathbb{R}$. The sum of the squared
Fourier coefficients of its real-valued version over sets outside the family consisting
of sets contained in the coordinates whose Boolean influence is at least $\epsilon^2
I^{-2}9^{-k}$, together with sets of size greater than $k$, equals the sum over sets of
size at most $k$ that are not contained in those coordinates; here $I$ is the total
influence of the real-valued version of $f$.
\end{theorem}

\begin{theorem}[Concentration on influential or large sets]
\lean{FABL.theorem9_28_spectrum_concentrated}
\proofsource{Analysis-of-Boolean-Functions-by-Ryan-ODonnell}{pdf-ffbe9c7b3306-p268-b003}
\label{FABL.theorem9_28_spectrum_concentrated}
\leanok
\uses{FABL.BooleanFunction, FABL.BooleanFunction.toReal, FABL.IsFourierSpectrumConcentratedOn, FABL.fourierWeightOutside_theorem9_28Family, FABL.lowDegreeOutsideInfluentialMass, FABL.lowDegreeOutsideInfluentialMass_le, FABL.sqrt_theorem9_28Threshold, FABL.theorem9_28Family, FABL.theorem9_28Threshold, FABL.totalInfluence}
\difficulty{3}
Let $f:\{-1,1\}^n\to\{-1,1\}$, let $\varepsilon>0$, and let $k$ be a nonnegative
integer. Suppose the total influence $I(f)$ is positive. Set
$\tau=\dfrac{\varepsilon^2}{I(f)^2}\left(\dfrac{1}{9}\right)^k$, and let $J$ consist of
the coordinates whose Boolean influence is at least $\tau$. Then the total squared
Fourier coefficient mass outside the family consisting of all subsets of $J$ together
with all subsets of $[n]$ of size greater than $k$ is at most $\varepsilon$.
\end{theorem}

\begin{definition}[Small theorem family]
\lean{FABL.theorem9_28SmallFamily}
\statementsource{Analysis-of-Boolean-Functions-by-Ryan-ODonnell}{pdf-ffbe9c7b3306-p269-b002}
\label{FABL.theorem9_28SmallFamily}
\leanok
\uses{FABL.BooleanFunction, FABL.influentialCoordinates, FABL.theorem9_28Threshold}
For a Boolean function $f$ on the $n$-dimensional Boolean cube, a real number
$\varepsilon$, and a nonnegative integer $k$, define the family of subsets
$S\subseteq\{1,\ldots,n\}$ that are contained in the coordinates whose Boolean influence
is at least $\varepsilon^2/I(f)^2\cdot 9^{-k}$ and satisfy $|S|\le k$, where $I(f)$ is
the total influence of $f$.
\end{definition}

\begin{theorem}[Small threshold-family exterior weight bound]
\lean{FABL.fourierWeightOutside_theorem9_28SmallFamily_le_add}
\label{FABL.fourierWeightOutside_theorem9_28SmallFamily_le_add}
\leanok
\uses{FABL.BooleanFunction, FABL.BooleanFunction.toReal, FABL.fourierWeightAbove, FABL.fourierWeightOutside, FABL.influentialCoordinates, FABL.theorem9_28Family, FABL.theorem9_28SmallFamily, FABL.theorem9_28Threshold}
\difficulty{4}
Let $n$ and $k$ be natural numbers, let $f$ be a sign-valued function on the
$n$-dimensional sign cube, and let $\epsilon\in\mathbb{R}$. The sum of the squared
Fourier coefficients of its real-valued version over sets outside the family of sets of
size at most $k$ contained in the coordinates whose Boolean influence is at least
$\epsilon^2 I^{-2}9^{-k}$ is at most the sum over sets outside the family consisting of
those coordinate subsets together with all sets of size greater than $k$, plus the sum
over sets of size greater than $k$; here $I$ is the total influence of the real-valued
version of $f$.
\end{theorem}

\begin{theorem}[Concentration on small influential sets]
\lean{FABL.theorem9_28_small_spectrum_concentrated}
\proofsource{Analysis-of-Boolean-Functions-by-Ryan-ODonnell}{
  pdf-ffbe9c7b3306-p268-b003,
  pdf-ffbe9c7b3306-p269-b001
}
\label{FABL.theorem9_28_small_spectrum_concentrated}
\leanok
\uses{FABL.BooleanFunction, FABL.BooleanFunction.toReal, FABL.IsFourierSpectrumConcentratedOn, FABL.IsFourierSpectrumConcentratedUpTo, FABL.fourierWeightAbove, FABL.fourierWeightAboveReal_natCast, FABL.fourierWeightOutside_theorem9_28SmallFamily_le_add, FABL.theorem9_28SmallFamily, FABL.theorem9_28_spectrum_concentrated, FABL.totalInfluence}
\difficulty{3}
Let $f:\{-1,1\}^n\to\{-1,1\}$, let $\varepsilon>0$, and let $k$ be a nonnegative
integer. Suppose the total influence $I(f)$ is positive and the total squared Fourier
coefficient mass on subsets of size greater than $k$ is at most $\varepsilon$. Set
$\tau=\dfrac{\varepsilon^2}{I(f)^2}\left(\dfrac{1}{9}\right)^k$, and let $J$ consist of
the coordinates whose Boolean influence is at least $\tau$. The total squared Fourier
coefficient mass outside the family of subsets of $J$ having size at most $k$ is at most
$2\varepsilon$.
\end{theorem}

\begin{definition}[Boolean coordinate projection]
\lean{FABL.booleanCoordinateProjection}
\statementsource{Analysis-of-Boolean-Functions-by-Ryan-ODonnell}{
  pdf-ffbe9c7b3306-p015-b012,
  pdf-ffbe9c7b3306-p212-b002,
  pdf-ffbe9c7b3306-p212-b004,
  pdf-ffbe9c7b3306-p269-b002
}
\label{FABL.booleanCoordinateProjection}
\leanok
\uses{FABL.BooleanFunction, FABL.BooleanFunction.toReal, FABL.Sign, FABL.projectOnCoordinates, FABL.uniformPMF}
For a Boolean function $f$ on the $n$-dimensional Boolean cube and a set of coordinates
$J\subseteq\{1,\ldots,n\}$, its coordinate projection onto $J$ is the real-valued
function obtained by projecting the real-valued version of $f$ onto those coordinates
with respect to the uniform probability measure on the sign set.
\end{definition}

\begin{definition}[Coordinate junta approximation]
\lean{FABL.coordinateJuntaApproximation}
\statementsource{Analysis-of-Boolean-Functions-by-Ryan-ODonnell}{pdf-ffbe9c7b3306-p269-b002}
\label{FABL.coordinateJuntaApproximation}
\leanok
\uses{FABL.BooleanFunction, FABL.booleanCoordinateProjection, FABL.thresholdSign}
For a Boolean function $f$ on the $n$-dimensional Boolean cube and a set of coordinates
$J\subseteq\{1,\ldots,n\}$, define its coordinate junta approximation by applying the
sign threshold map to the coordinate projection of $f$ onto $J$ at each input.
\end{definition}

\begin{theorem}[Fourier coefficients of coordinate projection]
\lean{FABL.fourierCoeff_booleanCoordinateProjection}
\proofsource{Analysis-of-Boolean-Functions-by-Ryan-ODonnell}{
  pdf-ffbe9c7b3306-p015-b012,
  pdf-ffbe9c7b3306-p212-b004,
  pdf-ffbe9c7b3306-p212-b005,
  pdf-ffbe9c7b3306-p213-b001
}
\label{FABL.fourierCoeff_booleanCoordinateProjection}
\leanok
\uses{FABL.BooleanFunction, FABL.BooleanFunction.toReal, FABL.FiniteFourierBasis.fourierCoeff_projectOnCoordinates, FABL.MultiIndex, FABL.Sign, FABL.booleanCoordinateProjection, FABL.fourierCoeff, FABL.projectOnCoordinates, FABL.signMultiIndexEquivFinset, FABL.signMultiIndexSupport, FABL.signSiteFourierBasis, FABL.signSiteFourierBasis_fourierCoeff_eq_fourierCoeff, FABL.uniformPMF}
\difficulty{3}
For every natural number $n$, every sign-valued function $f$ on the $n$-dimensional sign
cube, and coordinate sets $J,T\subseteq\{1,\ldots,n\}$, the Fourier coefficient at $T$
of the projection of $f$ onto coordinates in $J$ is the Fourier coefficient of its
real-valued version at $T$ if $T\subseteq J$, and is zero otherwise.
\end{theorem}

\begin{theorem}[Supremum norm of a sign-valued function]
\lean{FABL.productSupNorm_toReal_eq_one}
\label{FABL.productSupNorm_toReal_eq_one}
\leanok
\uses{FABL.BooleanFunction, FABL.BooleanFunction.toReal, FABL.abs_le_productSupNorm, FABL.productSupNorm}
\difficulty{4}
For every natural number $n$ and every function $f$ from the sign cube $\{-1,1\}^n$ to
$\{\pm1\}$, the supremum norm of its real-valued sign representation is $1$.
\end{theorem}

\begin{theorem}[Coordinate projection is bounded by one]
\lean{FABL.abs_booleanCoordinateProjection_le_one}
\label{FABL.abs_booleanCoordinateProjection_le_one}
\leanok
\uses{FABL.BooleanFunction, FABL.BooleanFunction.toReal, FABL.Sign, FABL.abs_projectOnCoordinates_le_productSupNorm, FABL.booleanCoordinateProjection, FABL.productSupNorm_toReal_eq_one, FABL.uniformPMF}
\difficulty{3}
Let $n$ be a natural number, let $f$ map the Boolean cube $\{-1,1\}^n$ to the sign set
$\{-1,1\}$, let $J$ be a subset of the coordinate set $\{1,\ldots,n\}$, and let
$x\in\{-1,1\}^n$. The absolute value of the coordinate projection of $f$ onto $J$,
evaluated at $x$, is at most $1$.
\end{theorem}

\begin{theorem}[Coordinate Junta Dependence]
\lean{FABL.coordinateJuntaApproximation_dependsOnly}
\proofsource{Analysis-of-Boolean-Functions-by-Ryan-ODonnell}{pdf-ffbe9c7b3306-p269-b002}
\label{FABL.coordinateJuntaApproximation_dependsOnly}
\leanok
\uses{FABL.BooleanFunction, FABL.BooleanFunction.toReal, FABL.DependsOnlyOnCoordinates, FABL.Sign, FABL.booleanCoordinateProjection, FABL.coordinateJuntaApproximation, FABL.projectOnCoordinates_dependsOnly, FABL.signValue, FABL.thresholdSign, FABL.uniformPMF}
\difficulty{3}
Let $n$ be a natural number, let $f:\{−1,1\}^n\to\{−1,1\}$, and let
$J\subseteq\{1,\ldots,n\}$ be a set of coordinates. Form a function by averaging $f$
over independent uniform choices of the coordinates outside $J$, while keeping the
coordinates in $J$ fixed, and then taking the sign of that average, with zero assigned
the positive sign. Its real-valued form depends only on the coordinates in $J$: any two
inputs agreeing on $J$ have the same value.
\end{theorem}

\begin{theorem}[Coordinate junta approximation is k junta]
\lean{FABL.coordinateJuntaApproximation_isKJunta}
\statementsource{Analysis-of-Boolean-Functions-by-Ryan-ODonnell}{
  pdf-ffbe9c7b3306-p044-b003,
  pdf-ffbe9c7b3306-p269-b002
}
\label{FABL.coordinateJuntaApproximation_isKJunta}
\leanok
\uses{FABL.BooleanFunction, FABL.BooleanFunction.toReal, FABL.IsKJunta, FABL.Sign, FABL.coordinateJuntaApproximation, FABL.projectOnCoordinates_dependsOnly, FABL.uniformPMF}
\difficulty{3}
% dataset-pipeline: stub -- prose pending build_informal_statements.py
Statement pending regeneration.
\end{theorem}

\begin{theorem}[Projection norm and Fourier weight]
\lean{FABL.uniformInner_booleanCoordinateProjection_self}
\label{FABL.uniformInner_booleanCoordinateProjection_self}
\leanok
\uses{FABL.BooleanFunction, FABL.BooleanFunction.toReal, FABL.booleanCoordinateProjection, FABL.fourierCoeff, FABL.fourierCoeff_booleanCoordinateProjection, FABL.fourierWeight, FABL.fourierWeightOutside, FABL.plancherel, FABL.sum_sq_fourierCoeff_eq_one}
\difficulty{5}
Let $f:\{−1,1\}^n\to\{−1,1\}$ be a Boolean function and let $J\subseteq\{1,\ldots,n\}$.
The squared uniform norm of its projection onto the coordinates in $J$ equals $1$ minus
the sum of the squared Fourier coefficients of $f$ indexed by subsets of
$\{1,\ldots,n\}$ that are not contained in $J$.
\end{theorem}

\begin{theorem}[Junta approximation inner product identity]
\lean{FABL.uniformInner_coordinateJuntaApproximation}
\label{FABL.uniformInner_coordinateJuntaApproximation}
\leanok
\uses{FABL.BooleanFunction, FABL.BooleanFunction.toReal, FABL.Sign, FABL.booleanCoordinateProjection, FABL.coordinateJuntaApproximation, FABL.coordinateJuntaApproximation_dependsOnly, FABL.productInner_uniformSign, FABL.projectOnCoordinates_eq_self_of_dependsOnly, FABL.projectOnCoordinates_selfAdjoint, FABL.signSiteFourierBasis, FABL.uniformPMF}
\difficulty{4}
Let $f:\{−1,1\}^n\to\{−1,1\}$ be a Boolean function and let $J\subseteq\{1,\ldots,n\}$.
The uniform inner product of $f$ with its sign-thresholded coordinate projection onto
$J$ equals the uniform inner product of that coordinate projection with the same
thresholded function.
\end{theorem}

\begin{theorem}[Projection and rounding inner product]
\lean{FABL.uniformInner_projection_rounding_eq_expect_abs}
\label{FABL.uniformInner_projection_rounding_eq_expect_abs}
\leanok
\uses{FABL.BooleanFunction, FABL.BooleanFunction.toReal, FABL.booleanCoordinateProjection, FABL.coordinateJuntaApproximation, FABL.signValue_thresholdSign, FABL.uniformInner}
\difficulty{4}
Let $f:\{−1,1\}^n\to\{−1,1\}$ and let $J\subseteq\{1,\ldots,n\}$. For each input $x$,
let $P_Jf(x)$ be the average of $f$ over inputs that agree with $x$ on $J$, and let
$R_Jf(x)$ be $1$ when $P_Jf(x)\geq0$ and $−1$ otherwise. Then the uniform inner product
of $P_Jf$ with $R_Jf$ equals the uniform average of $|P_Jf|$ over the cube.
\end{theorem}

\begin{theorem}[Projection self inner product bound]
\lean{FABL.uniformInner_projection_self_le_expect_abs}
\label{FABL.uniformInner_projection_self_le_expect_abs}
\leanok
\uses{FABL.BooleanFunction, FABL.abs_booleanCoordinateProjection_le_one, FABL.booleanCoordinateProjection, FABL.uniformInner}
\difficulty{3}
Let $f:\{−1,1\}^n\to\{−1,1\}$ and let $J\subseteq\{1,\ldots,n\}$. If $P_Jf(x)$ denotes
the average of $f$ over inputs agreeing with $x$ on $J$, then the uniform inner product
of $P_Jf$ with itself is at most the uniform average of $|P_Jf|$ over the cube.
\end{theorem}

\begin{theorem}[Two mul relative hamming dist coordinate junta approximation le]
\lean{FABL.two_mul_relativeHammingDist_coordinateJuntaApproximation_le}
\statementsource{Analysis-of-Boolean-Functions-by-Ryan-ODonnell}{pdf-ffbe9c7b3306-p269-b002}
\label{FABL.two_mul_relativeHammingDist_coordinateJuntaApproximation_le}
\leanok
\uses{FABL.BooleanFunction, FABL.BooleanFunction.toReal, FABL.coordinateJuntaApproximation, FABL.fourierWeightOutside, FABL.relativeHammingDist, FABL.uniformInner_booleanCoordinateProjection_self, FABL.uniformInner_coordinateJuntaApproximation, FABL.uniformInner_eq_one_sub_two_mul_relativeHammingDist, FABL.uniformInner_projection_rounding_eq_expect_abs, FABL.uniformInner_projection_self_le_expect_abs}
\difficulty{4}
% dataset-pipeline: stub -- prose pending build_informal_statements.py
Statement pending regeneration.
\end{theorem}

\begin{theorem}[Monotonicity of exterior Fourier weight]
\lean{FABL.fourierWeightOutside_antitone}
\label{FABL.fourierWeightOutside_antitone}
\leanok
\uses{FABL.fourierCoeff, FABL.fourierWeightOutside}
\difficulty{4}
Let $n$ be a natural number, let $f$ be a real-valued function on the $n$-dimensional
sign cube, and let $A$ and $B$ be families of subsets of the coordinates with
$A\subseteq B$. The sum of the squared Fourier coefficients of $f$ over sets outside $B$
is at most the corresponding sum over sets outside $A$.
\end{theorem}

\begin{theorem}[Theorem9 28 close to coordinate junta]
\lean{FABL.theorem9_28_close_to_coordinateJunta}
\statementsource{Analysis-of-Boolean-Functions-by-Ryan-ODonnell}{
  pdf-ffbe9c7b3306-p071-b007,
  pdf-ffbe9c7b3306-p269-b002
}
\label{FABL.theorem9_28_close_to_coordinateJunta}
\leanok
\uses{FABL.BooleanFunction, FABL.BooleanFunction.toReal, FABL.IsFourierSpectrumConcentratedUpTo, FABL.coordinateJuntaApproximation, FABL.fourierWeightOutside, FABL.fourierWeightOutside_antitone, FABL.influentialCoordinates, FABL.relativeHammingDist, FABL.theorem9_28SmallFamily, FABL.theorem9_28Threshold, FABL.theorem9_28_small_spectrum_concentrated, FABL.totalInfluence, FABL.two_mul_relativeHammingDist_coordinateJuntaApproximation_le}
\difficulty{4}
% dataset-pipeline: stub -- prose pending build_informal_statements.py
Statement pending regeneration.
\end{theorem}

\begin{theorem}[Approximation by a junta]
\lean{FABL.theorem9_28_exists_close_junta}
\proofsource{Analysis-of-Boolean-Functions-by-Ryan-ODonnell}{
  pdf-ffbe9c7b3306-p268-b003,
  pdf-ffbe9c7b3306-p269-b002
}
\label{FABL.theorem9_28_exists_close_junta}
\leanok
\uses{FABL.BooleanFunction, FABL.BooleanFunction.toReal, FABL.IsFourierSpectrumConcentratedUpTo, FABL.IsKJunta, FABL.coordinateJuntaApproximation, FABL.coordinateJuntaApproximation_isKJunta, FABL.influentialCoordinates, FABL.relativeHammingDist, FABL.theorem9_28Threshold, FABL.theorem9_28_close_to_coordinateJunta, FABL.totalInfluence}
\difficulty{1}
Let $f:\{-1,1\}^n\to\{-1,1\}$, let $\varepsilon>0$, and let $k$ be a nonnegative
integer. Suppose the total influence $I(f)$ is positive and the total squared Fourier
coefficient mass on subsets of size greater than $k$ is at most $\varepsilon$. Set
$\tau=\dfrac{\varepsilon^2}{I(f)^2}\left(\dfrac{1}{9}\right)^k$, and let $J$ consist of
the coordinates whose Boolean influence is at least $\tau$. Then there is a Boolean
function $g$ depending on at most $|J|$ coordinates such that the relative Hamming
distance between $f$ and $g$ is at most $\varepsilon$.
\end{theorem}

\begin{theorem}[Fourier coeff eq zero of total influence eq zero]
\lean{FABL.fourierCoeff_eq_zero_of_totalInfluence_eq_zero}
\statementsource{Analysis-of-Boolean-Functions-by-Ryan-ODonnell}{pdf-ffbe9c7b3306-p052-b002}
\label{FABL.fourierCoeff_eq_zero_of_totalInfluence_eq_zero}
\leanok
\uses{FABL.fourierCoeff, FABL.totalInfluence, FABL.totalInfluence_eq_sum_card_mul_sq_fourierCoeff}
\difficulty{4}
% dataset-pipeline: stub -- prose pending build_informal_statements.py
Statement pending regeneration.
\end{theorem}

\begin{theorem}[Zero total influence implies constancy]
\lean{FABL.eq_of_totalInfluence_toReal_eq_zero}
\label{FABL.eq_of_totalInfluence_toReal_eq_zero}
\leanok
\uses{FABL.BooleanFunction, FABL.BooleanFunction.toReal, FABL.fourierCoeff_eq_zero_of_totalInfluence_eq_zero, FABL.fourier_expansion, FABL.monomial, FABL.signValue_injective, FABL.totalInfluence}
\difficulty{4}
Let $n$ be a nonnegative integer and let $f:\{-1,1\}^n\to\{-1,1\}$. If the total
influence of the real-valued version of $f$ is zero, then $f$ takes the same value at
every pair of inputs; in particular, $f$ is constant.
\end{theorem}

\begin{theorem}[Zero total influence implies constancy]
\lean{FABL.isKJunta_zero_of_totalInfluence_eq_zero}
\label{FABL.isKJunta_zero_of_totalInfluence_eq_zero}
\leanok
\uses{FABL.BooleanFunction, FABL.BooleanFunction.toReal, FABL.IsKJunta, FABL.eq_of_totalInfluence_toReal_eq_zero, FABL.totalInfluence}
\difficulty{3}
Let $f:\{-1,1\}^n\to\{-1,1\}$. If the total influence of its real-valued sign
representation is zero, then $f$ depends on at most zero coordinates.
\end{theorem}

\begin{theorem}[Spectrum concentrated on empty of total influence eq zero]
\lean{FABL.spectrum_concentrated_on_empty_of_totalInfluence_eq_zero}
\statementsource{Analysis-of-Boolean-Functions-by-Ryan-ODonnell}{
  pdf-ffbe9c7b3306-p052-b002,
  pdf-ffbe9c7b3306-p069-b003
}
\label{FABL.spectrum_concentrated_on_empty_of_totalInfluence_eq_zero}
\leanok
\uses{FABL.BooleanFunction, FABL.BooleanFunction.toReal, FABL.IsFourierSpectrumConcentratedOn, FABL.fourierCoeff_eq_zero_of_totalInfluence_eq_zero, FABL.fourierWeight, FABL.fourierWeightOutside, FABL.totalInfluence}
\difficulty{4}
% dataset-pipeline: stub -- prose pending build_informal_statements.py
Statement pending regeneration.
\end{theorem}

\begin{definition}[Friedgut degree]
\lean{FABL.friedgutDegree}
\statementsource{Analysis-of-Boolean-Functions-by-Ryan-ODonnell}{pdf-ffbe9c7b3306-p269-b004}
\label{FABL.friedgutDegree}
\leanok
\uses{FABL.BooleanFunction, FABL.BooleanFunction.toReal, FABL.totalInfluence}
For a Boolean function $f$ on the $n$-dimensional sign cube and a real number
$\varepsilon$, define its Friedgut degree to be the natural-number ceiling of the total
influence of $f$, divided by $\varepsilon$.
\end{definition}

\begin{theorem}[Fourier concentration at the Friedgut degree]
\lean{FABL.spectrum_concentrated_up_to_friedgutDegree}
\proofsource{Analysis-of-Boolean-Functions-by-Ryan-ODonnell}{pdf-ffbe9c7b3306-p269-b004}
\label{FABL.spectrum_concentrated_up_to_friedgutDegree}
\leanok
\uses{FABL.BooleanFunction, FABL.BooleanFunction.toReal, FABL.IsFourierSpectrumConcentratedUpTo, FABL.IsFourierSpectrumConcentratedUpTo.mono_cutoff, FABL.friedgutDegree, FABL.isFourierSpectrumConcentratedUpTo_totalInfluence_div}
\difficulty{3}
Let $f:\{\pm1\}^n\to\{\pm1\}$ and let $\varepsilon>0$. Define the degree $d$ to be the
ceiling of the total influence of $f$ divided by $\varepsilon$. Then the sum of the
squared Fourier coefficients indexed by subsets of size greater than $d$ is at most
$\varepsilon$.
\end{theorem}

\begin{definition}[Friedgut coordinates]
\lean{FABL.friedgutCoordinates}
\statementsource{Analysis-of-Boolean-Functions-by-Ryan-ODonnell}{pdf-ffbe9c7b3306-p269-b004}
\label{FABL.friedgutCoordinates}
\leanok
\uses{FABL.BooleanFunction, FABL.friedgutDegree, FABL.influentialCoordinates, FABL.theorem9_28Threshold}
For a Boolean function $f$ on the $n$-dimensional sign cube and a real number
$\varepsilon$, define its Friedgut coordinate set to be the set of influential
coordinates of $f$ at the threshold supplied by Theorem 9.28 for $f$, $\varepsilon$, and
the Friedgut degree of $f$ at $\varepsilon$.
\end{definition}

\begin{definition}[Friedgut junta size bound]
\lean{FABL.friedgutJuntaSizeBound}
\statementsource{Analysis-of-Boolean-Functions-by-Ryan-ODonnell}{
  pdf-ffbe9c7b3306-p071-b007,
  pdf-ffbe9c7b3306-p269-b004
}
\label{FABL.friedgutJuntaSizeBound}
\leanok
\uses{FABL.BooleanFunction, FABL.BooleanFunction.toReal, FABL.friedgutDegree, FABL.totalInfluence}
For a Boolean function $f$ on the $n$-dimensional sign cube and a real number
$\varepsilon$, define the Friedgut junta size bound to be
$1+\frac{I(f)^3}{\varepsilon^2}9^{d}$, where $I(f)$ is the total influence of $f$ and
$d$ is its Friedgut degree at $\varepsilon$.
\end{definition}

\begin{theorem}[Card friedgut coordinates le]
\lean{FABL.card_friedgutCoordinates_le}
\statementsource{Analysis-of-Boolean-Functions-by-Ryan-ODonnell}{pdf-ffbe9c7b3306-p269-b004}
\label{FABL.card_friedgutCoordinates_le}
\leanok
\uses{FABL.BooleanFunction, FABL.BooleanFunction.toReal, FABL.card_theorem9_28_influentialCoordinates_le, FABL.friedgutCoordinates, FABL.friedgutDegree, FABL.totalInfluence}
\difficulty{1}
% dataset-pipeline: stub -- prose pending build_informal_statements.py
Statement pending regeneration.
\end{theorem}

\begin{theorem}[Friedgut spectrum concentrated]
\lean{FABL.friedgut_spectrum_concentrated}
\statementsource{Analysis-of-Boolean-Functions-by-Ryan-ODonnell}{
  pdf-ffbe9c7b3306-p269-b004,
  pdf-ffbe9c7b3306-p052-b002,
  pdf-ffbe9c7b3306-p069-b004
}
\label{FABL.friedgut_spectrum_concentrated}
\leanok
\uses{FABL.BooleanFunction, FABL.BooleanFunction.toReal, FABL.IsFourierSpectrumConcentratedOn, FABL.friedgutCoordinates, FABL.friedgutDegree, FABL.spectrum_concentrated_up_to_friedgutDegree, FABL.theorem9_28SmallFamily, FABL.theorem9_28_small_spectrum_concentrated, FABL.totalInfluence}
\difficulty{1}
% dataset-pipeline: stub -- prose pending build_informal_statements.py
Statement pending regeneration.
\end{theorem}

\begin{theorem}[Closeness of the coordinate junta approximation]
\lean{FABL.friedgut_coordinateJunta_close}
\proofsource{Analysis-of-Boolean-Functions-by-Ryan-ODonnell}{pdf-ffbe9c7b3306-p269-b004}
\label{FABL.friedgut_coordinateJunta_close}
\leanok
\uses{FABL.BooleanFunction, FABL.BooleanFunction.toReal, FABL.coordinateJuntaApproximation, FABL.fourierWeightOutside, FABL.fourierWeightOutside_antitone, FABL.friedgutCoordinates, FABL.friedgutDegree, FABL.influentialCoordinates, FABL.relativeHammingDist, FABL.spectrum_concentrated_up_to_friedgutDegree, FABL.theorem9_28SmallFamily, FABL.theorem9_28Threshold, FABL.theorem9_28_small_spectrum_concentrated, FABL.totalInfluence, FABL.two_mul_relativeHammingDist_coordinateJuntaApproximation_le}
\difficulty{4}
Let $f$ be a Boolean function on the $n$-dimensional sign cube, and suppose that
$\varepsilon>0$ and the total influence of $f$ is positive. The function obtained by
thresholding the coordinate projection of $f$ onto its Friedgut coordinate set at
parameter $\varepsilon$ has relative Hamming distance at most $\varepsilon$ from $f$.
\end{theorem}

\begin{theorem}[Friedgut junta approximation]
\lean{FABL.friedgut_junta}
\proofsource{Analysis-of-Boolean-Functions-by-Ryan-ODonnell}{pdf-ffbe9c7b3306-p269-b004}
\label{FABL.friedgut_junta}
\leanok
\uses{FABL.BooleanFunction, FABL.BooleanFunction.toReal, FABL.IsFourierSpectrumConcentratedOn, FABL.IsKJunta, FABL.card_friedgutCoordinates_le, FABL.coordinateJuntaApproximation, FABL.coordinateJuntaApproximation_isKJunta, FABL.friedgutCoordinates, FABL.friedgutDegree, FABL.friedgutJuntaSizeBound, FABL.friedgut_coordinateJunta_close, FABL.friedgut_spectrum_concentrated, FABL.isKJunta_zero_of_totalInfluence_eq_zero, FABL.relativeHammingDist, FABL.spectrum_concentrated_on_empty_of_totalInfluence_eq_zero, FABL.totalInfluence, FABL.totalInfluence_nonneg}
\difficulty{4}
Let $f$ be a Boolean function on the $n$-dimensional sign cube and let $\varepsilon>0$.
There exist a coordinate set $J$ and a Boolean function $g$ such that $|J|\le
1+I(f)^3\varepsilon^{-2}9^{d}$, where $I(f)$ is the total influence of $f$ and $d$ is
its Friedgut degree at parameter $\varepsilon$. The Fourier weight of $f$ outside the
frequencies $S\subseteq J$ with $|S|\le d$ is at most $2\varepsilon$; moreover, $g$
depends on at most $|J|$ coordinates and has relative Hamming distance at most
$\varepsilon$ from $f$.
\end{theorem}

\begin{definition}[DNF Friedgut degree]
\lean{FABL.dnfFriedgutDegree}
\statementsource{Analysis-of-Boolean-Functions-by-Ryan-ODonnell}{pdf-ffbe9c7b3306-p269-b005}
\label{FABL.dnfFriedgutDegree}
\leanok
\uses{FABL.dnfSwitchingDepth}
For a natural number $w$ and a real number $\varepsilon$, define the DNF Friedgut degree
to be the natural-number ceiling of $30w$ times the switching depth for parameter
$\varepsilon$.
\end{definition}

\begin{theorem}[Fourier concentration for bounded-width DNF]
\lean{FABL.dnf_spectrum_concentrated_up_to_dnfFriedgutDegree}
\proofsource{Analysis-of-Boolean-Functions-by-Ryan-ODonnell}{pdf-ffbe9c7b3306-p269-b005}
\label{FABL.dnf_spectrum_concentrated_up_to_dnfFriedgutDegree}
\leanok
\uses{FABL.BooleanFunction, FABL.BooleanFunction.toReal, FABL.HasDNFWidthLE, FABL.IsFourierSpectrumConcentratedUpTo, FABL.IsFourierSpectrumConcentratedUpTo.mono_cutoff, FABL.dnfFriedgutDegree, FABL.theorem4_22}
\difficulty{1}
Let $f$ be a sign-valued Boolean function on the $n$-dimensional Boolean cube that is
represented by a disjunctive normal form of width at most the natural number $w$. For
$0<\varepsilon\le 1$, the total squared Fourier weight of $f$, viewed as a real-valued
function, on sets of coordinates of size strictly greater than
$\lceil30w\lceil\log_2(3/\varepsilon)\rceil\rceil$ is at most $\varepsilon$.
\end{theorem}

\begin{definition}[DNF junta size bound]
\lean{FABL.dnfFriedgutJuntaSizeBound}
\statementsource{Analysis-of-Boolean-Functions-by-Ryan-ODonnell}{pdf-ffbe9c7b3306-p269-b006}
\label{FABL.dnfFriedgutJuntaSizeBound}
\leanok
\uses{FABL.dnfFriedgutDegree}
For a natural number $w$ and a real number $\varepsilon$, define the DNF junta size
bound to be $1+\frac{(2w)^3}{\varepsilon^2}9^d$, where $d$ is the DNF Friedgut degree
for $w$ and $\varepsilon$.
\end{definition}

\begin{theorem}[Exists close junta of has dnf width le]
\lean{FABL.exists_close_junta_of_hasDNFWidthLE}
\statementsource{Analysis-of-Boolean-Functions-by-Ryan-ODonnell}{
  pdf-ffbe9c7b3306-p269-b005,
  pdf-ffbe9c7b3306-p269-b006
}
\label{FABL.exists_close_junta_of_hasDNFWidthLE}
\leanok
\uses{FABL.BooleanFunction, FABL.BooleanFunction.toReal, FABL.HasDNFWidthLE, FABL.IsKJunta, FABL.card_theorem9_28_influentialCoordinates_le, FABL.coordinateJuntaApproximation, FABL.coordinateJuntaApproximation_isKJunta, FABL.dnfFriedgutDegree, FABL.dnfFriedgutJuntaSizeBound, FABL.dnf_spectrum_concentrated_up_to_dnfFriedgutDegree, FABL.influentialCoordinates, FABL.isKJunta_zero_of_totalInfluence_eq_zero, FABL.relativeHammingDist, FABL.theorem9_28Threshold, FABL.theorem9_28_close_to_coordinateJunta, FABL.totalInfluence, FABL.totalInfluence_le_two_mul_of_hasDNFWidthLE, FABL.totalInfluence_nonneg}
\difficulty{5}
% dataset-pipeline: stub -- prose pending build_informal_statements.py
Statement pending regeneration.
\end{theorem}
